\IfFileExists{stacks-project.cls}{%
\documentclass{stacks-project}
}{%
\documentclass{amsart}
}

% For dealing with references we use the comment environment
\usepackage{verbatim}
\newenvironment{reference}{\comment}{\endcomment}
%\newenvironment{reference}{}{}
\newenvironment{slogan}{\comment}{\endcomment}
\newenvironment{history}{\comment}{\endcomment}

% For commutative diagrams we use Xy-pic
\usepackage[all]{xy}

% We use 2cell for 2-commutative diagrams.
\xyoption{2cell}
\UseAllTwocells

% We use multicol for the list of chapters between chapters
\usepackage{multicol}

% This is generally recommended for better output
\usepackage{lmodern}
\usepackage[T1]{fontenc}

% For cross-file-references
\usepackage{xr-hyper}

% Package for hypertext links:
\usepackage{hyperref}

% For any local file, say "hello.tex" you want to link to please
% use \externaldocument[hello-]{hello}
\externaldocument[introduction-]{introduction}
\externaldocument[conventions-]{conventions}
\externaldocument[sets-]{sets}
\externaldocument[categories-]{categories}
\externaldocument[topology-]{topology}
\externaldocument[sheaves-]{sheaves}
\externaldocument[sites-]{sites}
\externaldocument[stacks-]{stacks}
\externaldocument[fields-]{fields}
\externaldocument[algebra-]{algebra}
\externaldocument[brauer-]{brauer}
\externaldocument[homology-]{homology}
\externaldocument[derived-]{derived}
\externaldocument[simplicial-]{simplicial}
\externaldocument[more-algebra-]{more-algebra}
\externaldocument[smoothing-]{smoothing}
\externaldocument[modules-]{modules}
\externaldocument[sites-modules-]{sites-modules}
\externaldocument[injectives-]{injectives}
\externaldocument[cohomology-]{cohomology}
\externaldocument[sites-cohomology-]{sites-cohomology}
\externaldocument[dga-]{dga}
\externaldocument[dpa-]{dpa}
\externaldocument[sdga-]{sdga}
\externaldocument[hypercovering-]{hypercovering}
\externaldocument[schemes-]{schemes}
\externaldocument[constructions-]{constructions}
\externaldocument[properties-]{properties}
\externaldocument[morphisms-]{morphisms}
\externaldocument[coherent-]{coherent}
\externaldocument[divisors-]{divisors}
\externaldocument[limits-]{limits}
\externaldocument[varieties-]{varieties}
\externaldocument[topologies-]{topologies}
\externaldocument[descent-]{descent}
\externaldocument[perfect-]{perfect}
\externaldocument[more-morphisms-]{more-morphisms}
\externaldocument[flat-]{flat}
\externaldocument[groupoids-]{groupoids}
\externaldocument[more-groupoids-]{more-groupoids}
\externaldocument[etale-]{etale}
\externaldocument[chow-]{chow}
\externaldocument[intersection-]{intersection}
\externaldocument[pic-]{pic}
\externaldocument[weil-]{weil}
\externaldocument[adequate-]{adequate}
\externaldocument[dualizing-]{dualizing}
\externaldocument[duality-]{duality}
\externaldocument[discriminant-]{discriminant}
\externaldocument[derham-]{derham}
\externaldocument[local-cohomology-]{local-cohomology}
\externaldocument[algebraization-]{algebraization}
\externaldocument[curves-]{curves}
\externaldocument[resolve-]{resolve}
\externaldocument[models-]{models}
\externaldocument[functors-]{functors}
\externaldocument[equiv-]{equiv}
\externaldocument[pione-]{pione}
\externaldocument[etale-cohomology-]{etale-cohomology}
\externaldocument[proetale-]{proetale}
\externaldocument[relative-cycles-]{relative-cycles}
\externaldocument[more-etale-]{more-etale}
\externaldocument[trace-]{trace}
\externaldocument[crystalline-]{crystalline}
\externaldocument[spaces-]{spaces}
\externaldocument[spaces-properties-]{spaces-properties}
\externaldocument[spaces-morphisms-]{spaces-morphisms}
\externaldocument[decent-spaces-]{decent-spaces}
\externaldocument[spaces-cohomology-]{spaces-cohomology}
\externaldocument[spaces-limits-]{spaces-limits}
\externaldocument[spaces-divisors-]{spaces-divisors}
\externaldocument[spaces-over-fields-]{spaces-over-fields}
\externaldocument[spaces-topologies-]{spaces-topologies}
\externaldocument[spaces-descent-]{spaces-descent}
\externaldocument[spaces-perfect-]{spaces-perfect}
\externaldocument[spaces-more-morphisms-]{spaces-more-morphisms}
\externaldocument[spaces-flat-]{spaces-flat}
\externaldocument[spaces-groupoids-]{spaces-groupoids}
\externaldocument[spaces-more-groupoids-]{spaces-more-groupoids}
\externaldocument[bootstrap-]{bootstrap}
\externaldocument[spaces-pushouts-]{spaces-pushouts}
\externaldocument[spaces-chow-]{spaces-chow}
\externaldocument[groupoids-quotients-]{groupoids-quotients}
\externaldocument[spaces-more-cohomology-]{spaces-more-cohomology}
\externaldocument[spaces-simplicial-]{spaces-simplicial}
\externaldocument[spaces-duality-]{spaces-duality}
\externaldocument[formal-spaces-]{formal-spaces}
\externaldocument[restricted-]{restricted}
\externaldocument[spaces-resolve-]{spaces-resolve}
\externaldocument[formal-defos-]{formal-defos}
\externaldocument[defos-]{defos}
\externaldocument[cotangent-]{cotangent}
\externaldocument[examples-defos-]{examples-defos}
\externaldocument[algebraic-]{algebraic}
\externaldocument[examples-stacks-]{examples-stacks}
\externaldocument[stacks-sheaves-]{stacks-sheaves}
\externaldocument[criteria-]{criteria}
\externaldocument[artin-]{artin}
\externaldocument[quot-]{quot}
\externaldocument[stacks-properties-]{stacks-properties}
\externaldocument[stacks-morphisms-]{stacks-morphisms}
\externaldocument[stacks-limits-]{stacks-limits}
\externaldocument[stacks-cohomology-]{stacks-cohomology}
\externaldocument[stacks-perfect-]{stacks-perfect}
\externaldocument[stacks-introduction-]{stacks-introduction}
\externaldocument[stacks-more-morphisms-]{stacks-more-morphisms}
\externaldocument[stacks-geometry-]{stacks-geometry}
\externaldocument[moduli-]{moduli}
\externaldocument[moduli-curves-]{moduli-curves}
\externaldocument[examples-]{examples}
\externaldocument[exercises-]{exercises}
\externaldocument[guide-]{guide}
\externaldocument[desirables-]{desirables}
\externaldocument[coding-]{coding}
\externaldocument[obsolete-]{obsolete}
\externaldocument[fdl-]{fdl}
\externaldocument[index-]{index}

% Theorem environments.
%
\theoremstyle{plain}
\newtheorem{theorem}[subsection]{Theorem}
\newtheorem{proposition}[subsection]{Proposition}
\newtheorem{lemma}[subsection]{Lemma}

\theoremstyle{definition}
\newtheorem{definition}[subsection]{Definition}
\newtheorem{example}[subsection]{Example}
\newtheorem{exercise}[subsection]{Exercise}
\newtheorem{situation}[subsection]{Situation}

\theoremstyle{remark}
\newtheorem{remark}[subsection]{Remark}
\newtheorem{remarks}[subsection]{Remarks}

\numberwithin{equation}{subsection}

% Macros
%
\def\lim{\mathop{\mathrm{lim}}\nolimits}
\def\colim{\mathop{\mathrm{colim}}\nolimits}
\def\Spec{\mathop{\mathrm{Spec}}}
\def\Hom{\mathop{\mathrm{Hom}}\nolimits}
\def\Ext{\mathop{\mathrm{Ext}}\nolimits}
\def\SheafHom{\mathop{\mathcal{H}\!\mathit{om}}\nolimits}
\def\SheafExt{\mathop{\mathcal{E}\!\mathit{xt}}\nolimits}
\def\Sch{\mathit{Sch}}
\def\Mor{\mathop{\mathrm{Mor}}\nolimits}
\def\Ob{\mathop{\mathrm{Ob}}\nolimits}
\def\Sh{\mathop{\mathit{Sh}}\nolimits}
\def\NL{\mathop{N\!L}\nolimits}
\def\CH{\mathop{\mathrm{CH}}\nolimits}
\def\proetale{{pro\text{-}\acute{e}tale}}
\def\etale{{\acute{e}tale}}
\def\QCoh{\mathit{QCoh}}
\def\Ker{\mathop{\mathrm{Ker}}}
\def\Im{\mathop{\mathrm{Im}}}
\def\Coker{\mathop{\mathrm{Coker}}}
\def\Coim{\mathop{\mathrm{Coim}}}

% Boxtimes
%
\DeclareMathSymbol{\boxtimes}{\mathbin}{AMSa}{"02}

%
% Macros for moduli stacks/spaces
%
\def\QCohstack{\mathcal{QC}\!\mathit{oh}}
\def\Cohstack{\mathcal{C}\!\mathit{oh}}
\def\Spacesstack{\mathcal{S}\!\mathit{paces}}
\def\Quotfunctor{\mathrm{Quot}}
\def\Hilbfunctor{\mathrm{Hilb}}
\def\Curvesstack{\mathcal{C}\!\mathit{urves}}
\def\Polarizedstack{\mathcal{P}\!\mathit{olarized}}
\def\Complexesstack{\mathcal{C}\!\mathit{omplexes}}
% \Pic is the operator that assigns to X its picard group, usage \Pic(X)
% \Picardstack_{X/B} denotes the Picard stack of X over B
% \Picardfunctor_{X/B} denotes the Picard functor of X over B
\def\Pic{\mathop{\mathrm{Pic}}\nolimits}
\def\Picardstack{\mathcal{P}\!\mathit{ic}}
\def\Picardfunctor{\mathrm{Pic}}
\def\Deformationcategory{\mathcal{D}\!\mathit{ef}}
\usepackage{cleveref}
\crefname{originallemma}{original lemma}{original lemmas}
\Crefname{originallemma}{Original lemma}{Original lemmas}
\makeatletter
\newlabel{spaces-properties-lemma-stalk-pullback}{{04K2}{source}{}{sourcepointer-1}{}}
\newlabel{spaces-properties-lemma-stalk-pullback@cref}{{[originallemma][0][]04K2}{[1][0][]source}{}{}{}}
\makeatother
\crefname{originallemma}{original lemma}{original lemmas}
\Crefname{originallemma}{Original lemma}{Original lemmas}
\makeatletter
\newlabel{spaces-cohomology-lemma-weak-chow}{{089J}{source}{}{sourcepointer-2}{}}
\newlabel{spaces-cohomology-lemma-weak-chow@cref}{{[originallemma][0][]089J}{[1][0][]source}{}{}{}}
\makeatother
\crefname{originallemma}{original lemma}{original lemmas}
\Crefname{originallemma}{Original lemma}{Original lemmas}
\makeatletter
\newlabel{topologies-lemma-surjective-proper-ph}{{0DES}{source}{}{sourcepointer-3}{}}
\newlabel{topologies-lemma-surjective-proper-ph@cref}{{[originallemma][0][]0DES}{[1][0][]source}{}{}{}}
\makeatother
\crefname{originallemma}{original lemma}{original lemmas}
\Crefname{originallemma}{Original lemma}{Original lemmas}
\makeatletter
\newlabel{etale-cohomology-lemma-describe-pullback-pi-ph}{{0DDW}{source}{}{sourcepointer-4}{}}
\newlabel{etale-cohomology-lemma-describe-pullback-pi-ph@cref}{{[originallemma][0][]0DDW}{[1][0][]source}{}{}{}}
\makeatother
\title{Injectivity of the unit F -> f\_* f\textasciicircum{}\{-1\}F for a surjective proper morphism, with image the equalizer}
\author{The Stacks Project Authors}
\date{stacks-project:spaces-more-cohomology.tex:local-snapshot-20261012; GNU FDL 1.2 or later}
\begin{document}
\maketitle
\noindent\textit{Editorial note.} The counted statement is the author's own lemma environment labelled lemma-surjective-proper (source lines 365-373); the proof environment beginning on the line after the following blank line (source line 375) is its own proof and closes at source line 403.\par\medskip
\noindent\textit{Editorial note.} The proof cites four results, all outside the chapter (Properties of Spaces, Cohomology of Spaces, Topologies, Etale Cohomology); each is declared as a tagged source pointer, and no step the counted argument needs is left to them: the four pointers are used exactly as the author uses them, as named inputs.\par\medskip
\noindent\textit{Editorial note.} The opening paragraph of the surrounding text ('The proper base change theorem for algebraic spaces follows from ...', source lines 360-363) is not part of the counted statement or proof and is not retained.\par\medskip
\noindent\textit{Editorial note.} The retained preamble is the chapter's own (sources/stacks-spaces-more-cohomology/preamble.tex, taken with its real line count), a verbatim excerpt of the native Xy-pic setup and theorem declarations; no macro was invented or redefined.\par\medskip
\noindent\textit{Editorial scope.} One original result of the Stacks Project's chapter Cohomology of Algebraic Spaces, counted once: for a surjective proper morphism f the unit F -> f\_*f\textasciicircum{}\{-1\}F is injective with image the equalizer.\par\medskip
\noindent\textit{Allowed standard prerequisites (editorial).}\begin{itemize}
\item Pullback of a sheaf along a morphism and its stalks. For a morphism f : Y -> X of algebraic spaces the unit F -> f\_*f\textasciicircum{}\{-1\}F, whose stalk at a geometric point y over x is F\_x via the canonical isomorphism (Properties of Spaces, declared as a tagged source pointer).. It computes the composite at a geometric point in the injectivity step (source lines 381-387).
\item Etale locality of the equalizer statement. The statement that the image is the equalizer is local on X for the etale topology, so one may assume Y is an affine scheme.. It is the reduction at source lines 390-391.
\item Chow's lemma for algebraic spaces. Every algebraic space Y admits a surjective proper morphism Z -> Y with Z a scheme (Cohomology of Spaces, declared as a tagged source pointer).. It produces the scheme Z at source lines 394-395.
\item Surjective proper morphisms of schemes are ph coverings, and the scheme case. A surjective proper morphism of schemes is a ph covering (Topologies), and for schemes the equalizer statement is known (Etale Cohomology, declared as tagged source pointers).. They transfer the result from Z to Y at source lines 397-402.
\end{itemize}\medskip
\ifdefined\hypertarget\hypertarget{sourcepointer-1}{}\fi
\noindent\textit{Source reference.} Original lemma 04K2: sources/stacks-chow/tags.txt, source lines 5958--5958. Cross-chapter pointer whose target carries the canonical tag 04K2 in the project tag list; declared so the reference resolves inside the assembled package, while every step the counted argument uses remains inside the retained spans.\par\medskip
\ifdefined\hypertarget\hypertarget{sourcepointer-2}{}\fi
\noindent\textit{Source reference.} Original lemma 089J: sources/stacks-chow/tags.txt, source lines 10499--10499. Cross-chapter pointer whose target carries the canonical tag 089J in the project tag list; declared so the reference resolves inside the assembled package, while every step the counted argument uses remains inside the retained spans.\par\medskip
\ifdefined\hypertarget\hypertarget{sourcepointer-3}{}\fi
\noindent\textit{Source reference.} Original lemma 0DES: sources/stacks-chow/tags.txt, source lines 16812--16812. Cross-chapter pointer whose target carries the canonical tag 0DES in the project tag list; declared so the reference resolves inside the assembled package, while every step the counted argument uses remains inside the retained spans.\par\medskip
\ifdefined\hypertarget\hypertarget{sourcepointer-4}{}\fi
\noindent\textit{Source reference.} Original lemma 0DDW: sources/stacks-chow/tags.txt, source lines 16781--16781. Cross-chapter pointer whose target carries the canonical tag 0DDW in the project tag list; declared so the reference resolves inside the assembled package, while every step the counted argument uses remains inside the retained spans.\par\medskip
\begin{lemma}
\label{lemma-surjective-proper}
Let $S$ be a scheme. Let $f : Y \to X$ be a surjective proper morphism
of algebraic spaces over $S$. Let $\mathcal{F}$ be a sheaf on $X_\etale$.
Then $\mathcal{F} \to f_*f^{-1}\mathcal{F}$ is injective with
image the equalizer of the two maps
$f_*f^{-1}\mathcal{F} \to g_*g^{-1}\mathcal{F}$ where
$g$ is the structure morphism $g : Y \times_X Y \to X$.
\end{lemma}

\begin{proof}
For any surjective morphism $f : Y \to X$ of algebraic spaces over $S$,
the map $\mathcal{F} \to f_*f^{-1}\mathcal{F}$ is injective.
Namely, if $\overline{x}$ is a geometric point of $X$, then we
choose a geometric point $\overline{y}$ of $Y$ lying over $\overline{x}$
and we consider
$$
\mathcal{F}_{\overline{x}} \to
(f_*f^{-1}\mathcal{F})_{\overline{x}} \to
(f^{-1}\mathcal{F})_{\overline{y}} = \mathcal{F}_{\overline{x}}
$$
See Properties of Spaces, Lemma \ref{spaces-properties-lemma-stalk-pullback}
for the last equality.

\medskip\noindent
The second statement is local on $X$ in the \'etale topology, hence we may
and do assume $Y$ is an affine scheme.

\medskip\noindent
Choose a surjective proper morphism $Z \to Y$ where $Z$ is a scheme, see
Cohomology of Spaces, Lemma \ref{spaces-cohomology-lemma-weak-chow}.
The result for $Z \to X$ implies the result for $Y \to X$.
Since $Z \to X$ is a surjective proper morphism of schemes
and hence a ph covering
(Topologies, Lemma \ref{topologies-lemma-surjective-proper-ph})
the result for $Z \to X$ follows from
\'Etale Cohomology, Lemma \ref{etale-cohomology-lemma-describe-pullback-pi-ph}
(in fact it is in some sense equivalent to this lemma).
\end{proof}

\end{document}
